% Part: applied-modal-logics
% Chapter: epistemic-logic
% Section: introduction

\documentclass[../../../include/open-logic-section]{subfiles}

\begin{document}

\olfileid[zh]{aml}{el}{int}

\olsection{引言}

正如模态逻辑处理\emph{模态命题}以及它们之间的衍推关系，认知逻辑处理\emph{认知命题}以及它们之间的衍推关系。我们不是把模态算子解释为表示可能性与必然性，而是以认知的或信念的方式解释这些一元联结词，用以刻画知识与信念。例如，我们可能想要表达像下面这样的断定：

\begin{enumerate}
\item 理查德知道卡尔加里在阿尔伯塔省。
\item 奥黛丽认为沙发上可能有只狗。
\item 理查德知道奥黛丽知道她的课在星期二。
\item 每个人都知道一年有 12 个月。
\end{enumerate}
当代认知逻辑常常追溯到\zhFullName{Jaakko Hintikka}{亚科·欣蒂卡}{欣蒂卡}的《知识与信念》（1962 年），而它写作的时代正是可能世界语义越来越频繁地用于逻辑学的时期。事实上，认知逻辑使用了大多数与其他模态逻辑相同的语义工具，但会对它们作不同的解释。主要的变化在于我们把\emph{可及关系}理解为表示什么。在认知逻辑中，它们表示某种形式的\emph{认知可能性}。我们将会看到，我们所刻画的这个认知概念会影响我们希望加诸可及关系之上的约束。我们还将看到对应理论在作认知解释时会变成什么样子。你会注意到，上面的例子提到了两个主体：理查德和奥黛丽，以及他们各自所知事物之间的关系。我们将要考察的认知逻辑将是多主体逻辑，在这样的逻辑中此类事情可以被表达。相反，单主体认知逻辑只会谈论某个个体知道或相信什么。

\end{document}
